\label{chap:taxonomia-monodromia}

The positional representation of a number records a sequence of symbols,
but does not necessarily preserve the state determining its continuation.
Triadic Modular Holography distinguishes the two levels. Nonadic phase indicates
the position of a reading within the nine-step cycle; the enriched
state additionally preserves orientation, Euclidean quotient, carry
transport, and boundary condition. Return of the phase does not therefore
imply return of the state.

Let \(S\) be a space of enriched states, let
\(T:S\to S\) be a transition, and let
\[
g:S\longrightarrow\mathbb Z/9\mathbb Z
\]
the phase, with \(g(Ts)=g(s)+1\). For a basis \(b\geq2\), a reader
\(d:S\to\{0,\ldots,b-1\}\) produces
\[
x(s)=\sum_{n\geq0}\frac{d(T^ns)}{b^{n+1}}.
\]
When a value admits the two usual positional expansions, we adopt the canonical representation that does not end in a constant tail \(b-1\). This convention distinguishes an ambiguity of notation from multiplicity proper to HMT, which arises from distinct enriched sections having the same evaluation.
The nonadic return operator is the restriction
\[
M=T^9:g^{-1}(r)\longrightarrow g^{-1}(r).
\]
The difference between \(T^9s=s\) and \(g(T^9s)=g(s)\) formalizes the distinction
between periodicity and monodromy.

The transition of the enriched state of TPK may be given by a relation rather than a
function. In that case, the preceding notation applies to the shift on
the space of infinite histories or to a deterministic realization whose state
has been enlarged with the information necessary to fix the continuation; it is not
presupposed that a visible block has a unique successor.

\begin{definition}[HMT Number and Arithmetic Presentation]
An HMT number is solely a compatible section
\(\boldsymbol x\) of the inverse system of surviving states. Given a
positional reader \(L\) whose domain contains \(\boldsymbol x\), the pair
\((\boldsymbol x,L)\) is a presentation or measurement of the number. The
Archimedean value published by that presentation is the singleton intersection of the
associated cylinders when their diameters tend to zero. Two sections with the
same value are not identified before applying the Archimedean quotient.
\end{definition}

\begin{definition}[Visible and enriched periodicities]
The output is eventually periodic when there exist \(N\geq0\) and \(p>0\)
such that
\[
d(T^{n+p}s)=d(T^ns),\qquad n\geq N.
\]
The state is eventually periodic when
\[
T^{N+p}s=T^Ns.
\]
The second condition implies the first. The converse requires that the reader
separate the tails of the orbit under consideration, that is, that
\[
 \bigl(d(T^{n+k}s)\bigr)_{k\geq0}
 =\bigl(d(T^{m+k}s)\bigr)_{k\geq0}
 \quad\Longrightarrow\quad T^ns=T^ms.
\]
\end{definition}

\begin{theorem}[Classification by phase return]
\label{thm:clasificacion-retorno-fase}
Let \(s\in S\).
\begin{enumerate}
  \item If the orbit of \(s\) reaches a state whose future output is
  constantly null, then \(x(s)\) has a terminating expansion and is
  rational.
  \item If the state is eventually periodic, then the output is
  eventually periodic and \(x(s)\in\mathbb Q\).
  \item If the reader separates the tails of the orbit and the orbit of \(s\) is not
  eventually periodic, then the output is not eventually periodic
  and \(x(s)\notin\mathbb Q\).
  \item Nonperiodicity does not distinguish between algebraic irrationality and
  transcendence. This distinction requires a polynomial certificate or an
  additional transcendence theorem.
\end{enumerate}
\end{theorem}

\begin{proof}
The first two parts are the classical characterization of the rationals through terminating or eventually periodic positional expansions. In the third, eventual periodicity of the output would produce two identical tails. The separation hypothesis would then identify the corresponding states, contradicting the nonperiodicity of the orbit. The fourth part follows from the existence of algebraic and transcendental irrationals with expansions that are not eventually periodic.
\end{proof}

\begin{corollary}[Nonadic Monodromy and Aperiodicity]
Suppose that the phase returns every nine steps, that the return orbit
\[
 \{M^ks:k\geq0\},\qquad M=T^9,
\]
is not eventually periodic and the reader separates the tails. Then the
output is not eventually periodic, and the number produced is irrational. In this
case, phase periodicity coexists with an aperiodic evolution of the
enriched state.
\end{corollary}

The hypothesis cannot be replaced by a single inequality \(M(s)\ne s\), nor
by a finite number of distinct returns. Relative to the initial state,
the declared transition and reader, the current window certifies nine
phase returns with change of state; it proves finite monodromy, but not
by itself the aperiodicity of the entire orbit.

\ifdefined\HMTRepetirResultadosTaxonomicos
\section{The local square of \texorpdfstring{\(e\)}{e} and the memory break}

For a decimal word \(a_1\cdots a_n\), we shall call a factorization a \emph{local square}
of length \(\ell\) at position \(i\)
\[
 a_i\cdots a_{i+2\ell-1}=vv,
 \qquad |v|=\ell.
\]
This notion belongs to combinatorics on words and must not be confused with
a period of the infinite tail.

\begin{theorem}[Finite characterization of the square \(1828^2\)]
Consider the \(243\) words of the N33 catalog and the frozen reader
\(w_6\mapsto w_{30}\mapsto d_1d_2d_3d_4\), with four triads in base one thousand.
There is a unique word whose reading contains a square of length four
beginning at the second decimal digit. It is
\[
 w_e=201101,
 \qquad
 718281828459=7\,\underbrace{1828\,1828}_{(1828)^2}\,459.
\]
The repetition is not prolonged to a third copy: the next required digit
would be \(1\), whereas the effective digit is \(4\). In the whole catalogue there
exists only one other square of length four,
\[
 575157518472=\underbrace{5751\,5751}_{\text{word square}}\,8472,
\]
but begins at the first digit.
\end{theorem}

\begin{proof}
The exhaustive enumeration evaluates the \(243\) words by exact
integer arithmetic. The square profiles for lengths \(1,\ldots,6\) contain,
respectively,
\[
240, 21, 3, 2, 0, 0
\]
occurrences, distributed among
\[
153, 19, 3, 2, 0, 0
\]
words. The two occurrences of length four are exactly those written
in the statement. Two independent integer enumerations reproduce the same
census, and their complete outputs agree term by term.
\end{proof}

The word \(w_e\) also has the lifting chain
\[
201101\mid121221\mid102011\mid012222\mid102011.
\]
The third and fifth visible blocks coincide. The states that carry them,
however, do not coincide. Before both emissions, their reductions
modulo \(27\) are
\[
(25,11,7,17,8,16),
\qquad
(7,8,4,23,26,7),
\]
and the subsequent Hensel quotients are
\[
(6,13,9,2,13,1),
\qquad
(15,0,3,19,23,25).
\]
Therefore, the return \(102011\) is a return of the visible projection, not
of the enriched state. This is the exact formulation of the intuition of a
local repetition that breaks: the quotient remembers a genealogy that the
residual block does not preserve.

The result belongs to the enriched reader already fixed by the system
APP--TRIT--TPK\@. The finite census exhaustively characterizes the appearance of the
local square; the coinductive prolongation of the section belongs to the
nonadic connection constructed in the
\cref{chap:numero-heisenberg-d108}. The combinatorial property of the prefix and
the generation of the complete expansion are therefore two levels
compatible with one and the same genealogy.

This formulation provides the exact framework for interpreting the
decimal beginning of \(e\),
\[
e=2.718281828459\ldots.
\]
The correct expression is no longer ``semiperiodicity'', but \emph{local word
square with visible return and memory rupture}.

\section{2 finite measures of the emitter \(468\to243\): frequency and structural priority}

There is no uniform distribution on all real numbers. There is, however,
a canonical measure on the finite catalog when the set of \(104\,976\) TPK seeds\@
is taken to be uniform. For a visible word \(w\), define
\[
\mu_{\mathrm{TPK}}(w)
=\frac{\#\{\text{seeds projecting onto }w\}}{104\,976}.
\]
The measure is not uniform. The \(243\) words have weights ranging
from \(144\) to \(2088\), with nine distinct multiplicity values.
The exact histogram, writing the weight as the base and the number of words
as the exponent, is
\[
144^{120},\ 288^{54},\ 432^{18},\ 576^9,\ 864^{15},\
1008^6,\ 1440^3,\ 1944^{12},\ 2088^6.
\]

For the three distinguished words one obtains
\[
\mu_{\mathrm{TPK}}(010211)=\frac{1008}{104\,976}=\frac7{729},
\]
\[
\mu_{\mathrm{TPK}}(201101)
=\mu_{\mathrm{TPK}}(121200)
=\frac{144}{104\,976}=\frac1{729}.
\]
The word associated with \(\pi\) has five preimages \(U_6\) and weight seven
times greater than the words associated with \(e\) and \(\varphi\). This result
proves a combinatorial hierarchy in the catalogue; by itself, it neither identifies
the words with the classical constants nor establishes a probability
on \(\mathbb R\).

Nor does it by itself define a total priority: six words have weight
\(1008\), whereas \(e\) and \(\varphi\) belong to the minimal class of
weight \(144\), shared by another \(118\) words. The HMT distinction must
be expressed by a structural profile, for example
\[
\mathcal P(w)=
\bigl(m_{\rm semilla}(w),m_{U_6}(w),r_{\rm supervivencia}(w),
s_{\rm frontera}(w),s_{\rm excepcional}(w)\bigr),
\]
with each coordinate defined before applying external labels. A scalar of ``priority'' further requires fixing, likewise in advance, how those coordinates are ordered or weighted.
\fi

\section{Classification of constants by genealogical antecedent, generating operation, and codomain}

The HMT classification separates properties that the ordinary representation
brings together under a single value:
\begin{center}
\begin{tabular}{p{0.25\textwidth}p{0.64\textwidth}}
\toprule
Class & Structural condition \\
\midrule
Definitive & The output reaches a null tail. \\
Rational periodic & The output is eventually periodic. \\
Rational monodromy & The output is periodic, but the enriched state
traverses a nontrivial fibre that the visible reader forgets. \\
Irracional coinductivo & The output is not eventually periodic and the
cylinders determine a unique value. \\
Recurrente local & A finite word contains returns or squares, without this imposing periodicity on the suffix. \\
Algebraic & The value satisfies a nonzero polynomial with integral coefficients; the
route must be accompanied by that certificate. \\
Transcendent & The value satisfies no nonzero integer polynomial; the
classification requires a transcendence theorem. \\
Distinguido HMT & The section satisfies a frozen structural predicate of
fiber, boundary, survival, or symmetry; it is a class transverse to the
preceding ones. \\
\bottomrule
\end{tabular}
\end{center}

Thus, HMT does not alter the correct arithmetic definitions of rationality,
algebraicity, and transcendence. It refines them by adding the genealogy of the state,
phase monodromy, and the multiplicity of the fibers that Archimedean
evaluation had forgotten.

In particular, \(\varphi\) is algebraic because it satisfies
\(X^2-X-1=0\), whereas the transcendence of \(e\) and \(\pi\) follows from
classical external theorems. These properties are compatible with the three
sections being distinguished by HMT\@. The dodecaphasic closure later constructs
the exact rational projection
\(\alpha_{12}=\pi_{12}(\alpha_{\mathrm{HMT}})\), and the compatible prolongation
constructs the coinductive section \(\alpha_{\mathrm{HMT}}\), in
\cref{chap:alpha-rev7}; neither is used in this taxonomy as an input
or published here as a second residence. The kernel \(E_{21/22}\)
produces a distinct analytic root, whose agreement with the integer route is
proved in the codomain of \(\operatorname{Tr}_9\). Section, rational
projection, analytic root, reader and value thus retain their own types.

Finally, an ``infinitesimal direction'' requires its own typing. For a finite reader \(r_m\), two distinct sections in the same fiber of \(r_m\) define a difference that is invisible at depth \(m\), although active in a subsequent prolongation. They do not thereby constitute a nonzero infinitesimal real. If two sections coincide under every faithful truncation of the inverse system, they are the same section. A genuine arithmetic infinitesimal would require the construction of an ordered or valued non-Archimedean extension and proof that the HMT operations descend to it.
